\chapter{Conclusão}
\label{cap:conclusao}

% -----------------------------------------------------------------
% Objetivo do capítulo:
% Fechar o trabalho retomando o problema, sintetizando o que foi
% entregue, reconhecendo os limites e indicando caminhos futuros.
% A conclusão não deve trazer resultados novos.
% -----------------------------------------------------------------


\section{Síntese da pesquisa}
\label{sec:sintese}

% Conteúdo a desenvolver:
%   - Retomar a pergunta de pesquisa e a proposição central.
%   - Resumir como a ferramenta foi construída para respondê-la.


\section{Contribuições efetivamente entregues}
\label{sec:contribuicoes_entregues}

% Conteúdo a desenvolver:
%   - Mapear cada contribuição prometida no Capítulo 1 ao
%     que foi de fato realizado.
%   - Destacar a arquitetura híbrida, a integração entre as
%     duas famílias de evidências e o mecanismo explicável.


\section{Limitações}
\label{sec:limitacoes}

% Conteúdo a desenvolver:
%   - Limitações de cobertura e qualidade de dados.
%   - Limitações da estratégia de avaliação.
%   - Reafirmar o caráter inferencial e não causal do sistema.


\section{Trabalhos futuros}
\label{sec:trabalhos_futuros}

% Conteúdo a desenvolver:
%   - Ampliar cobertura de sensores e de tipos de defeito.
%   - Incorporar realimentação de especialistas como fonte de
%     rótulo fraco.
%   - Generalizar a arquitetura para outros processos de
%     manufatura contínua com falhas multifatoriais.
%   - Direções identificadas durante a validação no caso
%     IBRAME 482530.
